% GENERATED FILE. Edit canonical lessons, then run npm run generate:books.
% canonical-source-hash: fnv1a64:b6ece98ec75b7597
% canonical-lessons: GE-C264-damals, GE-C264-einst, GE-C264-neulich, GE-C264-kuerzlich, GE-C264-vorgestern

\chapter{Damals und neulich}
\label{ch:ge-damals-und-neulich}
\hlchaptermodality{264}

\begin{chapteropening}
\textbf{By the end of this chapter:} \emph{I can place events in time: back then, long ago, the other day, recently and the day before yesterday.}

Complete the last lesson of chapter 264: I can place events in time: back then, long ago, the other day, recently and the day before yesterday.
\end{chapteropening}

\section[damals]{damals — back then, at that time}

\label{lesson:GE-C264-damals}



\begin{quote}
Before the new one: say the German for as long as, then the German for since then.
\end{quote}

\subsection*{You'll want to know: damals}
\textbf{damals} — \textquotedblleft{}back then, at that time\textquotedblright{}.

\begin{quote}
Damals hatten wir kein Handy.
\end{quote}

\begin{quote}
\emph{Back then we had no mobile phone.}
\end{quote}

For a time that is over. In a story it goes with the \emph{Präteritum}: \emph{Damals wohnte ich auf dem Land}.

\begin{cousinweb}
\emph{damals} is usually traced to Middle High German \emph{dō māles} \textquotedblleft{}at that time\textquotedblright{}, from \emph{dō} \textquotedblleft{}then\textquotedblright{} and \emph{Mal} \textquotedblleft{}time, occasion\textquotedblright{}; EtymDB lists no parent.
\end{cousinweb}

\begin{grammarlens}[title={What you've built}]
One more word for telling, explaining and describing at B1.
\end{grammarlens}

\subsection*{Your turn}
\begin{itemize}
\raggedright
  \item \emph{Say it:} \emph{damals}
  \item \emph{Say it:} \emph{damals}, once more
  \item \emph{Say it:} a sentence of your own with \emph{damals}
  \item \emph{Recall:} say the German for as long as, then the German for since then, then say \emph{damals} again
\end{itemize}

\begin{tcolorbox}[breakable,colback=teal!4,colframe=teal!35!black,title={Before you move on}]
What does damals mean? (\textquotedblleft{}Back then, at that time\textquotedblright{}.) Say it once more.
\end{tcolorbox}
\section[einst]{einst — once, long ago; one day}

\label{lesson:GE-C264-einst}



\begin{quote}
Before the new one: say the German for since then, then the German for back then.
\end{quote}

\subsection*{You'll want to know: einst}
\textbf{einst} — \textquotedblleft{}once, long ago; one day\textquotedblright{}.

\begin{quote}
Hier stand einst eine Burg.
\end{quote}

\begin{quote}
\emph{Long ago a castle stood here.}
\end{quote}

Literary. It can also look ahead: \emph{einst wirst du es verstehen}.

\begin{cousinweb}
\emph{einst} is inherited from Middle High German \emph{einest}, from Old High German \emph{einēst} \textquotedblleft{}once\textquotedblright{}, built on \emph{ein} \textquotedblleft{}one\textquotedblright{} (EtymDB).
\end{cousinweb}

\begin{grammarlens}[title={What you've built}]
One more word for telling, explaining and describing at B1.
\end{grammarlens}

\subsection*{Your turn}
\begin{itemize}
\raggedright
  \item \emph{Say it:} \emph{einst}
  \item \emph{Say it:} \emph{einst}, once more
  \item \emph{Say it:} a sentence of your own with \emph{einst}
  \item \emph{Recall:} say the German for since then, then the German for back then, then say \emph{einst} again
\end{itemize}

\begin{tcolorbox}[breakable,colback=teal!4,colframe=teal!35!black,title={Before you move on}]
What does einst mean? (\textquotedblleft{}Once, long ago; one day\textquotedblright{}.) Say it once more.
\end{tcolorbox}
\section[neulich]{neulich — the other day, recently}

\label{lesson:GE-C264-neulich}



\begin{quote}
Before the new one: say the German for back then, then the German for once.
\end{quote}

\subsection*{You'll want to know: neulich}
\textbf{neulich} — \textquotedblleft{}the other day, recently\textquotedblright{}.

\begin{quote}
Neulich habe ich meine alte Freundin getroffen.
\end{quote}

\begin{quote}
\emph{The other day I met my old friend.}
\end{quote}

Conversational; for one event a short time ago, usually with the \emph{Perfekt}.

\begin{cousinweb}
\emph{neulich} is built on \emph{neu} \textquotedblleft{}new\textquotedblright{}, an inherited Germanic word, with the ending \emph{-lich}: newly, lately. EtymDB lists no parent for the adverb itself.
\end{cousinweb}

\begin{grammarlens}[title={What you've built}]
One more word for telling, explaining and describing at B1.
\end{grammarlens}

\subsection*{Your turn}
\begin{itemize}
\raggedright
  \item \emph{Say it:} \emph{neulich}
  \item \emph{Say it:} \emph{neulich}, once more
  \item \emph{Say it:} a sentence of your own with \emph{neulich}
  \item \emph{Recall:} say the German for back then, then the German for once, then say \emph{neulich} again
\end{itemize}

\begin{tcolorbox}[breakable,colback=teal!4,colframe=teal!35!black,title={Before you move on}]
What does neulich mean? (\textquotedblleft{}The other day, recently\textquotedblright{}.) Say it once more.
\end{tcolorbox}
\section[kürzlich]{kürzlich — recently, lately}

\label{lesson:GE-C264-kuerzlich}



\begin{quote}
Before the new one: say the German for once, then the German for the other day.
\end{quote}

\subsection*{You'll want to know: kürzlich}
\textbf{kürzlich} — \textquotedblleft{}recently, lately\textquotedblright{}.

\begin{quote}
Die Bäckerei hat kürzlich zugemacht.
\end{quote}

\begin{quote}
\emph{The bakery closed recently.}
\end{quote}

A little more formal than \emph{neulich}; common in written reports.

\begin{cousinweb}
\emph{kürzlich} is built on \emph{kurz} \textquotedblleft{}short\textquotedblright{}, which reached Old High German \emph{kurz} from Latin \emph{curtus} \textquotedblleft{}shortened\textquotedblright{} (EtymDB): a short time ago.
\end{cousinweb}

\begin{grammarlens}[title={What you've built}]
One more word for telling, explaining and describing at B1.
\end{grammarlens}

\subsection*{Your turn}
\begin{itemize}
\raggedright
  \item \emph{Say it:} \emph{kürzlich}
  \item \emph{Say it:} \emph{kürzlich}, once more
  \item \emph{Say it:} a sentence of your own with \emph{kürzlich}
  \item \emph{Recall:} say the German for once, then the German for the other day, then say \emph{kürzlich} again
\end{itemize}

\begin{tcolorbox}[breakable,colback=teal!4,colframe=teal!35!black,title={Before you move on}]
What does kürzlich mean? (\textquotedblleft{}Recently, lately\textquotedblright{}.) Say it once more.
\end{tcolorbox}
\section[vorgestern]{vorgestern — the day before yesterday}

\label{lesson:GE-C264-vorgestern}



\begin{quote}
Before the new one: say the German for the other day, then the German for recently.
\end{quote}

\subsection*{You'll want to know: vorgestern}
\textbf{vorgestern} — \textquotedblleft{}the day before yesterday\textquotedblright{}.

\begin{quote}
Vorgestern war ich krank, aber heute geht es mir gut.
\end{quote}

\begin{quote}
\emph{The day before yesterday I was ill, but today I am well.}
\end{quote}

The partner of \emph{übermorgen}, the day after tomorrow.

\begin{cousinweb}
\emph{vorgestern} joins \emph{vor} \textquotedblleft{}before\textquotedblright{} and \emph{gestern} \textquotedblleft{}yesterday\textquotedblright{}, an inherited Germanic word and a cousin of the \emph{yester-} in English \emph{yesterday}; EtymDB lists no parent for the compound.
\end{cousinweb}

\begin{grammarlens}[title={What you've built}]
One more word for telling, explaining and describing at B1.
\end{grammarlens}

\subsection*{Your turn}
\begin{itemize}
\raggedright
  \item \emph{Say it:} \emph{vorgestern}
  \item \emph{Say it:} \emph{vorgestern}, once more
  \item \emph{Say it:} a sentence of your own with \emph{vorgestern}
  \item \emph{Recall:} say the German for the other day, then the German for recently, then say \emph{vorgestern} again
\end{itemize}

\begin{tcolorbox}[breakable,colback=teal!4,colframe=teal!35!black,title={Before you move on}]
What does vorgestern mean? (\textquotedblleft{}The day before yesterday\textquotedblright{}.) Say it once more.
\end{tcolorbox}
